\chapter*{Introdução}
\addcontentsline{toc}{chapter}{Introdução}
\markboth{INTRODUÇÃO}{INTRODUÇÃO}

%%%%%%%%%%%%%%%%%%%%%%%%%%%%%%%%%%%%%%%%%%%%%%%%%%%%%%%%%%%%%%%%%%%%%%%%%%%%%%

\section{Sobre este livro}

%mbxSTARTIGNORE
\VolOneIntroExtrahead
%mbxENDIGNORE
% pretext (mbx) is always the two volume set
%mbxlatex
%mbxlatex \subsection{Sobre o Volume I}
%mbxlatex 

Este primeiro volume destina-se a um curso semestral de análise básica.
Com o segundo volume, o curso passa a ter duração de um ano.
O livro teve origem nas minhas notas de aula para a disciplina Math 444,
ministrada na University of Illinois at Urbana-Champaign (UIUC), no semestre
de outono de 2009. Acrescentei o capítulo sobre espaços métricos para a
disciplina Math 521, na University of Wisconsin--Madison (UW)\@.
O Volume II foi incluído para a disciplina Math 4143/4153, na Oklahoma State
University (OSU)\@.
Em geral, essas disciplinas têm como pré-requisito um curso básico de
demonstrações, que pode usar, por exemplo, \cite{Hammack}, \cite{GIAM} ou
\cite{DW}.

O livro pode ser usado tanto em um curso básico, destinado a estudantes que
talvez não pretendam fazer pós-graduação (como a disciplina UIUC 444), quanto
em um curso semestral mais avançado, que também aborde tópicos como espaços
métricos (como a disciplina UW 521).
Seguem algumas sugestões de organização para um curso semestral.
Um curso mais lento, como UIUC 444:
\begin{center}
\S0.3, \S1.1--\S1.4, \S2.1--\S2.5, \S3.1--\S3.4, \S4.1--\S4.2,
\S5.1--\S5.3, \S6.1--\S6.3
\end{center}
Um curso mais rigoroso, que aborde espaços métricos e avance
consideravelmente mais rápido (por exemplo, UW 521):
\begin{center}
\S0.3, \S1.1--\S1.4, \S2.1--\S2.5, \S3.1--\S3.4, \S4.1--\S4.2,
\S5.1--\S5.3, \S6.1--\S6.2, \S7.1--\S7.6
\end{center}
Também é possível oferecer um curso mais rápido, sem espaços métricos, que
abranja todas as seções dos capítulos 0 a 6. O número aproximado de aulas
indicado nas notas das seções até o Capítulo 6 é uma estimativa bastante
grosseira, elaborada para o curso mais lento.
Os primeiros capítulos do livro também podem ser usados em um curso
introdutório de demonstrações, como ocorre, por exemplo, na disciplina Math
201 da Iowa State University, em que este livro é usado junto com
\emph{Book of Proof}, de Hammack~\cite{Hammack}.

Com o Volume II\@, é possível oferecer um curso anual que aborde tópicos de
várias variáveis. Nesse caso, pode fazer sentido cobrir a maior parte do
primeiro volume no primeiro semestre e deixar os espaços métricos para o
início do segundo semestre.

O início do Volume I segue, em certa medida, o programa tradicional da
disciplina Math 444 da UIUC e, por isso, tem algumas semelhanças com
\emph{Introduction to Real Analysis}, de Bartle e Sherbert \cite{BS}, que era
o livro adotado na UIUC\@.
Uma diferença importante é que definimos a integral de Riemann por meio de
somas de Darboux, e não de partições etiquetadas. A abordagem de Darboux é muito
mais adequada a um curso deste nível.

Nossa abordagem permite concluir, em um semestre, um curso como UIUC 444 e,
ainda assim, dedicar algum tempo à troca de limites e terminar com o teorema
de Picard sobre existência e unicidade de soluções de equações diferenciais
ordinárias.
Esse teorema é um excelente exemplo de aplicação de muitos resultados
demonstrados no livro. Para estudantes mais avançados, o conteúdo pode ser
abordado mais rapidamente, de modo que se chegue aos espaços métricos e se
demonstre o teorema de Picard usando o teorema do ponto fixo, como é usual.

Há outros livros excelentes. Meu favorito é
\emph{Principles of Mathematical Analysis}, de Rudin~\cite{Rudin:baby},
conhecido com carinho como \emph{baby Rudin} (para distingui-lo de seu outro
excelente livro de análise, chamado \emph{big Rudin}). Inspirei-me bastante
nas ideias de Rudin. No entanto, esse livro é um pouco mais avançado e
ambicioso do que o curso apresentado aqui.
Para quem pretende continuar estudando matemática, o livro de Rudin é um
ótimo investimento. Uma alternativa mais barata e um pouco mais simples é
\emph{Introduction to Analysis}, de Rosenlicht~\cite{Rosenlicht}.
Há também \emph{Introduction to Real Analysis}, de William Trench, disponível
gratuitamente para download~\cite{Trench}.

\medskip

Uma observação sobre o estilo de algumas demonstrações: prefiro apresentar
diretamente ou pela contrapositiva muitas demonstrações que tradicionalmente
são feitas por contradição. Embora o livro inclua demonstrações por
contradição, recorro a esse método apenas quando a contrapositiva parece
muito pouco natural ou quando a contradição decorre rapidamente.
A contradição tem maior probabilidade de confundir estudantes iniciantes,
pois nela falamos de objetos que não existem.

Procuro evitar formalismos desnecessários quando não são úteis.
Além disso, à medida que avançamos no livro, as demonstrações e a linguagem
tornam-se um pouco menos formais, pois omitimos cada vez mais detalhes para
evitar sobrecarregar o texto.

Como regra geral, uso $\coloneqq$ em vez de $=$ para definir um objeto, e não
simplesmente para indicar uma igualdade. Uso esse símbolo com mais frequência
do que é usual, para dar ênfase.
Uso-o até mesmo quando o contexto é \myquote{local,} isto é, posso definir
uma função $f(x) \coloneqq x^2$ apenas para um exercício ou exemplo.

\medskip

Por fim, gostaria de agradecer a Jana Ma\v{r}\'ikov\'a,
Glen Pugh, Paul Vojta, Frank Beatrous, S\"{o}nmez \c{S}ahuto\u{g}lu,
Jim Brandt, Kenji Kozai, Arthur Busch,  Anton Petrunin,
Mark Meilstrup, Harold P.\ Boas, Atilla Y{\i}lmaz,
Thomas Mahoney, Scott Armstrong e Paul Sacks,
Matthias Weber, Manuele Santoprete,
Robert Niemeyer e Amanullah Nabavi,
por adotarem o livro em suas disciplinas e me darem muitas sugestões úteis.
Frank Beatrous escreveu as extensões da versão da University of Pittsburgh,
que inspiraram muitas das adições ao livro.
Também gostaria de agradecer a
Dan Stoneham, Jeremy Sutter, Eliya Gwetta, Daniel Pimentel-Alarc\'on,
Steve Hoerning, Yi Zhang, Nicole Caviris,
Kristopher Lee, Baoyue Bi, Hannah Lund,
Trevor Mannella, Mitchel Meyer, Gregory Beauregard,
Chase Meadors, Andreas Giannopoulos, Nick Nelsen,
Ru Wang, Trevor Fancher, Brandon Tague, Wang KP\@,
Wai Yan Pong, Sam Merat, Judah Nouriyelian,
Arnold Cross, Jesse Wallace,
Adnan Hashem Mohamed,
Nikita Borisov, Bob Strain,
Salven V.\ DeMartino,
Xuechi Wang,
Hugo Santarem de Araujo,
Geoffrey Sired,
Karabo Mathope (0x03B4 no GitHub),
Samyak Jain,
um ou dois leitores anônimos e, de modo geral, todos os estudantes de minhas
turmas
pelas sugestões e por identificarem erros e falhas de digitação.
%mbxlatex       Durante a elaboração deste livro, o autor recebeu apoio parcial
%mbxlatex       dos auxílios NSF DMS-0900885 e DMS-1362337.

% if in the two volume set then this will expand to the intro for vol II
%mbxSTARTIGNORE
\VolTwoIntro
%mbxENDIGNORE
% pretext (mbx) is always the two volume set
%mbxlatex
%mbxlatex \subsection{Sobre o Volume II}
%mbxlatex 
%mbxlatex \input{frag-vol2-intro.tex}
%mbxlatex 

%%%%%%%%%%%%%%%%%%%%%%%%%%%%%%%%%%%%%%%%%%%%%%%%%%%%%%%%%%%%%%%%%%%%%%%%%%%%%%

\sectionnewpage
\section{Sobre análise} \label{sec:aboutra}

Análise é o ramo da matemática que trata de desigualdades e limites. Este
curso aborda os conceitos mais básicos da análise. Seu objetivo é familiarizar
o leitor com demonstrações rigorosas em análise e também estabelecer uma base
sólida para o cálculo de uma variável (e de várias variáveis, caso o Volume II
também seja considerado).

O cálculo ensinou você, estudante, a usar a matemática sem explicar por que
aquilo que aprendeu é verdadeiro. Para usar ou ensinar matemática com
eficácia, não basta saber \emph{o que} é verdadeiro; é preciso saber
\emph{por que} é verdadeiro. Este curso mostra \emph{por que} o cálculo é
válido. Seu propósito é proporcionar uma boa compreensão dos conceitos de
limite, derivada e integral.

Vamos recorrer a uma analogia.
Um mecânico que aprendeu a trocar o óleo, consertar faróis quebrados e carregar
a bateria, mas não entende como o motor funciona, só será capaz de realizar
essas tarefas simples. Não conseguirá diagnosticar e resolver problemas por
conta própria.
Um professor do ensino médio que não compreende a definição da integral de
Riemann ou da derivada talvez não consiga responder adequadamente a todas as
perguntas dos estudantes.
Até hoje me lembro de várias afirmações sem sentido que ouvi do meu professor
de cálculo no ensino médio: ele simplesmente não entendia o conceito de
limite, embora soubesse \myquote{fazer} os exercícios do livro.

\medskip

Começamos discutindo o sistema dos números reais e, sobretudo, sua propriedade
de completude, que serve de base para tudo o que vem a seguir.
Depois, tratamos da forma mais simples de limite: o limite de uma sequência.
Em seguida, estudamos funções de uma variável, continuidade e derivada.
Definimos então a integral de Riemann e demonstramos o teorema fundamental
do cálculo. Também discutimos sequências de funções e a troca de limites.
Por fim, apresentamos uma introdução aos espaços métricos.

\medskip

Vejamos a diferença mais importante entre análise e álgebra.
Em álgebra, demonstramos igualdades diretamente: provamos que um objeto,
talvez um número, é igual a outro.
Em análise, geralmente demonstramos desigualdades e, para isso, fazemos
estimativas.
Para ilustrar essa ideia, considere a seguinte afirmação.

\medskip

\emph{Seja $x$ um número real. Se $x < \epsilon$ vale para todo número real
$\epsilon > 0$, então $x \leq 0$.}

\medskip

Essa afirmação resume, em linhas gerais, o que fazemos em análise.
Suponha que queiramos demonstrar a igualdade $x = 0$. Em análise,
demonstramos duas desigualdades: $x \leq 0$ e $x \geq 0$.
Para demonstrar $x \leq 0$, mostramos que $x < \epsilon$ para todo
$\epsilon > 0$.
Para demonstrar $x \geq 0$, mostramos que $x > -\epsilon$ para todo
$\epsilon > 0$.

\medskip

O termo \emph{análise real} é um nome um tanto inadequado. Prefiro usar
simplesmente \emph{análise}. A outra área, a \emph{análise complexa}, na
verdade se apoia neste material, em vez de constituir algo completamente
distinto. Além disso, um curso mais avançado de análise real trataria com
frequência dos números complexos.
Acredito que essa nomenclatura seja um resquício histórico.

\medskip

Vamos ao que interessa\ldots


%%%%%%%%%%%%%%%%%%%%%%%%%%%%%%%%%%%%%%%%%%%%%%%%%%%%%%%%%%%%%%%%%%%%%%%%%%%%%%

\sectionnewpage
\section{Teoria básica dos conjuntos} \label{sec:basicset}
\index{teoria dos conjuntos}

%mbxINTROSUBSECTION

\sectionnotes{1--3 aulas (parte do material pode ser omitida, abordada
brevemente ou deixada como leitura)}

Antes de começar a falar de análise, precisamos fixar alguns termos.
A análise moderna\footnote{O termo \myquote{moderna} refere-se ao período
que vai do fim do século XIX até o presente.} usa a linguagem dos conjuntos;
por isso, é por aí que começamos.
Falamos de conjuntos de maneira relativamente informal, usando o que costuma
ser chamado de \myquote{\myindex{teoria ingênua dos conjuntos}.} Não se
preocupe: essa é a abordagem adotada pela maioria dos matemáticos, e é difícil
ter problemas com ela.
Espera-se que o leitor já tenha visto os fundamentos da teoria dos conjuntos
e da escrita de demonstrações; esta seção deve servir como uma breve
recapitulação.

\subsection{Conjuntos}

\begin{defn}
Um \emph{\myindex{conjunto}} é uma coleção de objetos chamados
\emph{elementos}\index{elemento} ou \emph{membros}\index{membro}. Um conjunto
sem elementos é chamado de \emph{\myindex{conjunto vazio}} e é denotado por
$\emptyset$\glsadd{not:emptyset} (ou, às vezes, por $\{ \}$).
\end{defn}

Imagine um conjunto como um clube com determinados associados. Por exemplo,
os estudantes que jogam xadrez são membros do clube de xadrez.
O mesmo estudante pode pertencer a vários clubes diferentes.
No entanto, não leve a analogia longe demais. Um conjunto é definido apenas
pelos membros que o formam; dois conjuntos com os mesmos membros são o mesmo
conjunto.

Na maior parte do tempo, consideraremos conjuntos de números. Por exemplo, o
conjunto
\glsadd{not:set}
\begin{equation*}
S \coloneqq \{ 0, 1, 2 \}
\end{equation*}
contém os três elementos 0, 1 e 2.
Com o símbolo \myquote{$\coloneqq$},\glsadd{not:def} indicamos que estamos
definindo o que é $S$, em vez de apenas mostrar uma igualdade.
Escrevemos\glsadd{not:eltof}
\begin{equation*}
1 \in S
\end{equation*}
para indicar que o número 1 pertence ao conjunto $S$. Em outras palavras, 1
é um elemento de $S$.
Às vezes, queremos dizer que dois elementos pertencem a um conjunto $S$; por
isso, usamos \myquote{$1,2 \in S$} como abreviação de
\myquote{$1 \in S$ e $2 \in S$.}

De modo semelhante, escrevemos\glsadd{not:noteltof}
\begin{equation*}
7 \notin S
\end{equation*}
para indicar que o número 7 não pertence a $S$, isto é, que 7 não é membro
de $S$.

Os elementos de todos os conjuntos considerados pertencem a algum conjunto
que chamamos de \emph{\myindex{universo}}. Para simplificar, muitas vezes
tomamos como universo o conjunto que contém apenas os elementos de nosso
interesse.
Em geral, o universo fica subentendido pelo contexto e não é mencionado
explicitamente. Neste curso, muitas vezes o universo será o conjunto dos
números reais.

Embora os elementos de um conjunto sejam frequentemente números, outros
objetos, como conjuntos, também podem ser elementos de um conjunto.
Além disso, dois conjuntos podem ter alguns elementos em comum. Por exemplo,
\begin{equation*}
T \coloneqq \{ 0, 2 \}
\end{equation*}
contém os números 0 e 2. Nesse caso, todos os elementos de $T$ também
pertencem a $S$. Escrevemos $T \subset S$.\glsadd{not:subset}
Veja o diagrama da \figureref{fig:sets}.
%More formally we make the
%following definition.

\begin{myfigureht}
\myincludepdft{sets}{%
Diagrama dos conjuntos de exemplo S e T. O conjunto S contém 0, 1 e 2,
e o conjunto T contém 0 e 2. O elemento 7 não pertence a nenhum dos
conjuntos.}
\caption{Diagrama dos conjuntos de exemplo $S$ e de seu subconjunto $T$.\label{fig:sets}}
\end{myfigureht}

\begin{defn}
%FIXMEevillayouthack
\pagebreak[2]
\leavevmode
\begin{enumerate}[(i)]
\item
Um conjunto $A$ é um \emph{\myindex{subconjunto}} de um conjunto $B$ se
$x \in A$ implica $x \in B$, e escrevemos $A \subset B$.
Isto é, todos os elementos de $A$ também são elementos de $B$. Às vezes,
escrevemos $B \supset A$ com o mesmo significado.
\item
Dois conjuntos $A$ e $B$ são \emph{\myindex{iguais}}\glsadd{not:equal} se
$A \subset B$ e $B \subset A$. Escrevemos $A = B$.
Isto é, $A$ e $B$ contêm exatamente os mesmos elementos.
Se $A$ e $B$ não são iguais, escrevemos $A \neq B$.
\item
Um conjunto $A$ é um \emph{\myindex{subconjunto pr\'oprio}} de $B$ se
$A \subset B$ e $A \neq B$. Escrevemos $A \subsetneq B$.\glsadd{not:subsetneq}
\end{enumerate}
\end{defn}

No exemplo dos conjuntos $S$ e $T$ definidos acima, $T \subset S$, mas
$T \neq S$. Portanto, $T$ é um
subconjunto próprio de $S$.
Se $A = B$, então $A$ e $B$ são apenas dois nomes para exatamente o mesmo
conjunto.

Para definir novos conjuntos, costuma-se usar a
\emph{\myindex{notação de constru\c{c}\~ao de conjuntos}},\glsadd{not:setbuilder}
\begin{equation*}
\bigl\{ x \in A : P(x) \bigr\} .
\end{equation*}
Essa notação representa um subconjunto de $A$ formado por todos os elementos
de $A$ que satisfazem a propriedade $P(x)$.
Usando $S = \{ 0, 1, 2 \}$ como acima, $\{ x \in S : x \neq 2 \}$
é o conjunto $\{ 0, 1 \}$.
Às vezes, a notação é abreviada como $\bigl\{ x : P(x) \bigr\}$; nesse caso,
o conjunto $A$ não é mencionado quando fica subentendido pelo contexto.
Além disso, a condição $x \in A$ pode ser substituída por uma fórmula para
tornar a notação mais fácil de ler.

\begin{example}
Os conjuntos a seguir usam as notações usuais.
\begin{enumerate}[(i)]
\item O conjunto dos \emph{\myindex{n\'umeros naturais}},
\glsadd{not:natnums}
$\N \coloneqq \{ 1, 2, 3, \ldots
\}$.
\item O conjunto dos \emph{\myindex{n\'umeros inteiros}},
\glsadd{not:ints} 
$\Z \coloneqq \{ 0, -1, 1, -2, 2, \ldots
\}$.
\item O conjunto dos \emph{\myindex{n\'umeros racionais}},
\glsadd{not:rationals} 
$\Q \coloneqq \bigl\{ \frac{m}{n} :  m, n \in \Z \text{ e } n \neq 0 \bigr\}$.
\item O conjunto dos números naturais pares, $\{  2m : m \in \N \}$.
\item O conjunto dos números reais,
\glsadd{not:reals} 
$\R$.
\end{enumerate}

Observe que $\N \subset \Z \subset \Q \subset \R$.
\end{example}

Formamos novos conjuntos a partir de conjuntos já existentes por meio de
algumas operações naturais.

\begin{defn}
\pagebreak[2]
\leavevmode
\begin{enumerate}[(i)]
\item
A \emph{\myindex{uni\~ao}} de dois conjuntos $A$ e $B$ é definida por
\glsadd{not:union}
\begin{equation*}
A \cup B \coloneqq \{ x : x \in A \text{ ou } x \in B \} .
\end{equation*}
\item
A \emph{\myindex{interse\c{c}\~ao}} de dois conjuntos $A$ e $B$ é definida por
\glsadd{not:intersection}
\begin{equation*}
A \cap B \coloneqq \{ x : x \in A \text{ e } x \in B \} .
\end{equation*}
\item
O \emph{complemento de $B$ em relação a $A$\index{complemento em rela\c{c}\~ao a}}
(ou a \emph{\myindex{diferen\c{c}a de conjuntos}} de $A$ e $B$) é definido por
\glsadd{not:setminus}
\begin{equation*}
A \setminus B \coloneqq \{ x : x \in A \text{ e } x \notin B \} .
\end{equation*}
\item
Chamamos de \emph{\myindex{complemento}} de $B$ e escrevemos $B^c$
\glsadd{not:complement} em vez de $A \setminus B$ quando o conjunto $A$ é
todo o universo ou o conjunto que obviamente contém $B$, conforme fique
subentendido pelo contexto.
\item
Dois conjuntos $A$ e $B$ são \emph{\myindex{disjuntos}} se $A \cap B =
\emptyset$.
\end{enumerate}
\end{defn}

A notação $B^c$ talvez pareça um pouco vaga neste momento. Se $B$ é um
subconjunto dos números reais $\R$, então $B^c$ significa
$\R \setminus B$. Se $B$ é naturalmente um subconjunto dos números naturais,
então $B^c$ é $\N \setminus B$. Se houver risco de ambiguidade, usamos a
notação de diferença de conjuntos $A \setminus B$.

\begin{myfigureht}
\myincludepdft{figsetop}{%
Diagramas de Venn das operações com conjuntos: primeiro, um diagrama da união
de A e B, com a região interna de A ou de B sombreada.
Depois, um diagrama da interseção de A e B, em que somente a interseção está
sombreada. Em seguida, um diagrama da diferença de A e B, em que tudo o que
está em A, mas não em B, está sombreado. Por fim, um diagrama do complemento
de B, em que tudo o que não está em B está sombreado.}
\caption{Diagramas de Venn das operações com conjuntos; a região sombreada
representa o resultado da operação.\label{figsetop}}
\end{myfigureht}
Ilustramos essas operações com os
\emph{diagramas de Venn}\index{diagrama de Venn} na \figureref{figsetop}.
Vamos agora demonstrar um dos teoremas mais básicos sobre conjuntos e lógica.

\begin{thm}[De Morgan]\index{teorema de De Morgan}
Sejam $A, B, C$ conjuntos. Então
\begin{equation*}
{(B \cup C)}^c = B^c \cap C^c , \qquad
{(B \cap C)}^c = B^c \cup C^c ,
\end{equation*}
ou, mais geralmente,
\begin{equation*}
A \setminus (B \cup C) = (A \setminus B) \cap (A \setminus C) , 
\qquad
A \setminus (B \cap C) = (A \setminus B) \cup (A \setminus C) .
\end{equation*}
\end{thm}

\begin{proof}
A primeira igualdade decorre da segunda se considerarmos o conjunto $A$ como
nosso \myquote{universo.}

Vamos demonstrar que
$A \setminus (B \cup C) = (A \setminus B) \cap (A \setminus C)$.
Para isso, lembre-se da definição de igualdade entre conjuntos. Primeiro,
precisamos mostrar que, se $x \in A \setminus (B \cup C)$, então
$x \in (A \setminus B) \cap (A \setminus C)$. Também precisamos mostrar que,
se $x \in (A \setminus B) \cap (A \setminus C)$, então
$x \in A \setminus (B \cup C)$.
Suponhamos, então, que $x \in A \setminus (B \cup C)$. Nesse caso, $x$
pertence a $A$, mas não pertence a $B$ nem a $C$. Portanto, $x$ pertence a
$A$ e não pertence a $B$, isto é, $x \in A \setminus B$. De modo semelhante,
$x \in A \setminus C$. Logo,
$x \in (A \setminus B) \cap (A \setminus C)$.
Por outro lado, suponha que
$x \in (A \setminus B) \cap (A \setminus C)$. Em particular,
$x \in (A \setminus B)$, de modo que $x \in A$ e $x \notin B$. Além disso, como
$x \in (A \setminus C)$, temos $x \notin C$.
Portanto, $x \in A \setminus (B \cup C)$.

A demonstração da outra igualdade fica como exercício.
\end{proof}

Chamamos o resultado acima de \emph{Teorema}; a maioria dos resultados recebe
o nome de \emph{Proposição}, e alguns são chamados de \emph{Lema} (um
resultado que conduz a outro) ou \emph{Corolário} (uma consequência rápida do
resultado anterior).
Não atribua importância excessiva a essa nomenclatura. Parte dela é
tradicional, e parte decorre de escolhas de estilo.
Não é necessariamente verdade que um \emph{Teorema} seja sempre
\myquote{mais importante} do que uma \emph{Proposição} ou um \emph{Lema}.

\medskip

Também precisaremos intersectar ou unir vários conjuntos ao mesmo tempo. Se
houver apenas um número finito deles, basta aplicar repetidamente a operação
de união ou de interseção. Suponha, porém, que tenhamos uma coleção infinita
de conjuntos (um conjunto de conjuntos)
$\{ A_1, A_2, A_3, \ldots \}$. Definimos
\glsadd{not:infunion}
\glsadd{not:infintersection}
\begin{align*}
& \bigcup_{n=1}^\infty A_n \coloneqq \{ x : x \in A_n \text{ para algum } n \in \N
\} , \\
& \bigcap_{n=1}^\infty A_n \coloneqq \{ x : x \in A_n \text{ para todo } n \in \N
\} .
\end{align*}

Também podemos ter conjuntos indexados por dois números naturais. Por
exemplo, podemos ter o conjunto de conjuntos
$\{ A_{1,1}, A_{1,2}, A_{2,1}, A_{1,3}, A_{2,2}, A_{3,1}, \ldots \}$. Então
escrevemos
\begin{equation*}
\bigcup_{n=1}^\infty \bigcup_{m=1}^\infty A_{n,m}
=
\bigcup_{n=1}^\infty \left( \bigcup_{m=1}^\infty A_{n,m} \right) .
\end{equation*}
De modo análogo, procedemos com as interseções.

Não é difícil ver que podemos tomar as uniões em qualquer ordem. No entanto,
em geral, não podemos trocar a ordem de uniões e interseções sem uma
demonstração.
Por exemplo,
\begin{equation*}
\bigcup_{n=1}^\infty
\bigcap_{m=1}^\infty
\{ k \in \N : mk < n \}
=
\bigcup_{n=1}^\infty \emptyset = \emptyset .
\end{equation*}
No entanto,
\begin{equation*}
\bigcap_{m=1}^\infty
\bigcup_{n=1}^\infty
\{ k \in \N : mk < n \}
=
\bigcap_{m=1}^\infty
\N
=
\N.
\end{equation*}

Às vezes, o conjunto de índices não é o conjunto dos números naturais. Nesse
caso, precisamos de uma notação mais geral. Suponha que $I$ seja um conjunto
e que, para cada $\lambda \in I$, exista um conjunto $A_\lambda$. Então
definimos
\glsadd{not:arbitunion}
\glsadd{not:arbitintersection}
\begin{equation*}
\bigcup_{\lambda \in I} A_\lambda \coloneqq \{ x : x \in A_\lambda \text{ para algum }
\lambda \in I
\} ,
\qquad
\bigcap_{\lambda \in I} A_\lambda \coloneqq \{ x : x \in A_\lambda \text{ para todo }
\lambda \in I
\} .
\end{equation*}

\subsection{Indu\c{c}\~ao}

Quando uma afirmação envolve um número natural arbitrário,
um método de demonstração comum é o princípio da indução.
Começamos com o conjunto dos números naturais $\N = \{ 1,2,3,\ldots \}$ e
consideramos sua ordem usual,
isto é, $1 < 2 < 3 < 4 < \cdots$.
Dizer que $S \subset \N$ possui um \emph{\myindex{menor elemento}} significa
que existe $x \in S$ tal que, para todo
$y \in S$, temos $x \leq y$.

Os números naturais $\N$, ordenados da maneira usual,
satisfazem a \emph{\myindex{propriedade da boa ordena\c{c}\~ao}}.
Adotamos essa propriedade como um \myindex{axioma}; simplesmente a supomos
verdadeira.

%mbxSTARTIGNORE
\theoremstyle{plain}
\newtheorem*{wellordprop}{Propriedade da boa ordena\c{c}\~ao de $\N$}
\hypertarget{wop:link}{}%
\begin{wellordprop}
Todo subconjunto não vazio de $\N$ possui um elemento mínimo (ou seja, um
menor elemento).
\end{wellordprop}
%mbxENDIGNORE
%mbx <axiom xml:id="wop_link" number="">
%mbx   <title>Propriedade da boa ordenação de <m>\N</m></title>
%mbx   <statement><p>
%mbx   Todo subconjunto não vazio de <m>\N</m> possui um elemento mínimo (ou seja, um menor elemento)
%mbx   </p></statement>
%mbx </axiom>

O \emph{princ\'ipio da \myindex{indu\c{c}\~ao}}\index{princ\'ipio da indu\c{c}\~ao}
\'e o teorema a seguir, que em certo sentido\footnote{Para sermos
completamente rigorosos, essa equival\^encia s\'o vale se tamb\'em
admitirmos como axioma que $n-1$ existe para todo n\'umero natural maior que
$1$, o que fazemos. Neste livro, supomos v\'alidas todas as regras usuais
da aritm\'etica.}
\'e equivalente \`a propriedade da boa ordena\c{c}\~ao dos n\'umeros naturais.

\begin{thm}[Princ\'ipio da indu\c{c}\~ao] \label{induction:thm}
Seja $P(n)$ uma afirma\c{c}\~ao que depende de um n\'umero natural~$n$.
Suponha que
\begin{enumerate}[(i)]
\item \emph{(\myindex{base da indu\c{c}\~ao})} $P(1)$ é verdadeira.
\item \emph{(\myindex{passo indutivo})} Se $P(n)$ é verdadeira, então
$P(n+1)$ também é verdadeira.
\end{enumerate}
Então $P(n)$ é verdadeira para todo $n \in \N$.
\end{thm}

\begin{proof}
Seja $S$ o conjunto dos números naturais $n$ para os quais $P(n)$ não é
verdadeira. Suponha, por contradição, que $S$ não seja vazio.
Pela propriedade da boa ordenação, $S$ possui um menor elemento.
Seja $m \in S$ o menor elemento de $S$. Pela hipótese, $1 \notin S$. Logo,
$m > 1$, e $m-1$ também é um número natural.
Como $m$ é o menor elemento de $S$, sabemos que $P(m-1)$ é verdadeira.
Porém, o passo indutivo afirma que $P(m-1+1) = P(m)$ é verdadeira, o
que contradiz a afirmação de que $m \in S$. Portanto, o conjunto $S$ é vazio, e
$P(n)$ é verdadeira para todo $n \in \N$.
\end{proof}

A suposição de que
$P(n)$ é verdadeira na afirmação \myquote{se $P(n)$ \'e verdadeira,
então $P(n+1)$ \'e verdadeira}
costuma ser chamada de \emph{\myindex{hip\'otese de indu\c{c}\~ao}}.
Às vezes, é conveniente iniciar a indução em um número diferente
de 1; nesse caso, isso altera apenas a numeração.
Por exemplo, se começarmos em $n=0$, então $P(0)$ é a base da
indução, e o passo indutivo deve ser demonstrado para todo
$n \in \N \cup \{ 0 \}$.
Concluímos, então, que $P(n)$ é verdadeira para todo
$n \in \N \cup \{ 0 \}$.

\begin{example}
Vamos provar que, para todo $n \in \N$,
\begin{equation*}
2^{n-1} \leq n! \qquad \text{(lembre-se de que } n! = 1 \cdot 2 \cdot 3 \cdots n\text{)}.
\end{equation*}
Seja $P(n)$ a afirmação de que
$2^{n-1} \leq n!$ é verdadeira.
Tomando $n=1$, vemos que $P(1)$ é verdadeira.

Suponha que $P(n)$ seja verdadeira, isto é, que
$2^{n-1} \leq n!$ valha. Multiplicando ambos os lados por 2, obtemos
\begin{equation*}
2^n \leq 2(n!) .
\end{equation*}
Como $2 \leq (n+1)$ para $n \in \N$, temos
$2(n!) \leq (n+1)(n!) = (n+1)!$. Isto é,
\begin{equation*}
2^n \leq 2(n!) \leq  (n+1)!,
\end{equation*}
e, assim, $P(n+1)$ é verdadeira. Pelo princípio da indução,
$P(n)$ é verdadeira para todo $n \in \N$. Em outras palavras,
$2^{n-1} \leq n!$ é verdadeira para todo $n \in \N$.
\end{example}

\begin{example} \label{example:geometricsum}
Afirmamos que, para todo $c \neq 1$
e todo $n \in \N \cup \{ 0 \}$,
\begin{equation*}
1 + c + c^2 + \cdots + c^n = \frac{1-c^{n+1}}{1-c} .
\end{equation*}

Demonstração: Fixemos $c \neq 1$.
É fácil verificar que a igualdade vale para $n=0$.
Suponha que ela valha para algum $n$ (possivelmente $n=0$).
Então
\begin{equation*}
\begin{split}
1 + c + c^2 + \cdots + c^n + c^{n+1} & =
( 1 + c + c^2 + \cdots + c^n ) + c^{n+1} \\
& = \frac{1-c^{n+1}}{1-c}  + c^{n+1} \\
& = \frac{1-c^{n+1}  + (1-c)c^{n+1}}{1-c} \\
& = \frac{1-c^{n+2}}{1-c} .
\end{split}
\end{equation*}
Pelo princípio da indução, a igualdade vale para todo
$n \in \N \cup \{ 0 \}$.
\end{example}

Às vezes, é mais fácil usar no passo indutivo o fato de que
$P(k)$ é verdadeira para todo $k = 1,2,\ldots,n$, e não apenas para $k=n$.
Esse princípio é chamado de \emph{\myindex{indu\c{c}\~ao forte}} e é
equivalente ao princípio de indução usual apresentado acima. A demonstração
dessa equivalência fica como exercício.

\begin{thm}[Princ\'ipio da indu\c{c}\~ao forte]
\index{princ\'ipio da indu\c{c}\~ao forte}\index{indu\c{c}\~ao forte}
Seja $P(n)$ uma afirmação que depende de um número natural $n$. Suponha que
\begin{enumerate}[(i)]
\item \emph{(\myindex{base da indu\c{c}\~ao})} $P(1)$ é verdadeira.
\item \emph{(\myindex{passo indutivo})} Se $P(k)$ é verdadeira para todo
$k = 1,2,\ldots,n$, então $P(n+1)$ é verdadeira.
\end{enumerate}
Então $P(n)$ é verdadeira para todo $n \in \N$.
\end{thm}

\subsection{Fun\c{c}\~oes}

Informalmente,
uma \emph{\myindex{fun\c{c}\~ao no sentido da teoria dos conjuntos}}\index{fun\c{c}\~ao}
$f$, que leva um conjunto $A$ a um conjunto $B$,
\'e uma aplica\c{c}\~ao que associa a cada $x \in A$ um \'unico $y \in B$.
Escrevemos $f \colon A \to B$.\glsadd{not:func}
Uma fun\c{c}\~ao $f \colon A \to B$ é como uma caixa-preta na qual
inserimos um elemento de $A$; a fun\c{c}\~ao nos devolve um elemento de $B$.
Às vezes, $f$ tamb\'em é chamada de \emph{\myindex{aplica\c{c}\~ao}} ou de
\emph{\myindex{mapa}}, e dizemos que $f$ \emph{leva $A$ em $B$}.

Uma fun\c{c}\~ao de exemplo
$f \colon S \to T$, de $S \coloneqq \{ 0, 1, 2 \}$ em $T \coloneqq \{ 0, 2 \}$,
pode ser definida atribuindo-se $f(0) \coloneqq 2$, $f(1) \coloneqq 2$ e
$f(2) \coloneqq 0$.
%
Muitas vezes, as fun\c{c}\~oes s\~ao definidas por algum tipo de
f\'ormula; no entanto, conv\'em pensar em uma fun\c{c}\~ao simplesmente como
uma tabela muito grande de valores.
A sutileza é que uma mesma fun\c{c}\~ao pode ter v\'arias f\'ormulas,
todas definindo a mesma fun\c{c}\~ao.
Para muitas fun\c{c}\~oes, n\~ao existe uma f\'ormula que exprima seus valores.
Para definir rigorosamente a no\c{c}\~ao de fun\c{c}\~ao, precisamos primeiro
definir o produto cartesiano.

\begin{defn}
Sejam $A$ e $B$ conjuntos. O \emph{\myindex{produto cartesiano}}
é o conjunto de tuplas definido por
\glsadd{not:cartprod}
\begin{equation*}
A \times B \coloneqq
\bigl\{ (x,y) : x \in A, y \in B \bigr\} .
\end{equation*}
\end{defn}

Por exemplo, $\{ a,b \} \times \{ c , d\} =
\bigl\{
(a,c), (a,d), (b,c), (b,d)
\bigr\}$.
Um exemplo mais elaborado é o conjunto $[0,1] \times [0,1]$: um
subconjunto do plano delimitado por um quadrado com v\'ertices $(0,0)$,
$(0,1)$, $(1,0)$ e $(1,1)$.
Quando $A$ e $B$ s\~ao o mesmo conjunto, costuma-se usar o expoente 2 para
denotar esse produto. Por exemplo, $[0,1]^2 =
[0,1] \times [0,1]$ ou $\R^2 = \R \times \R$ (o plano cartesiano).

\begin{defn}
Uma \emph{fun\c{c}\~ao} $f \colon A \to B$ é um subconjunto $f$ de $A \times B$
tal que, para cada $x \in A$, existe um \'unico $y \in B$ para o qual
$(x,y) \in f$.
Escrevemos $f(x) = y$.
Se considerarmos $f$ simplesmente como um subconjunto de $A \times B$,
chamamos esse conjunto de
\emph{\myindex{gr\'afico}} da fun\c{c}\~ao.

O conjunto $A$ é chamado de \emph{\myindex{dom\'inio}} de $f$ (e,
às vezes, é denotado de modo confuso por $D(f)$). O conjunto
\begin{equation*}
R(f) \coloneqq \{ y \in B : \text{existe um } x \in A \text{ tal que }
f(x)=y \}
\end{equation*}
Esse conjunto é chamado de \emph{\myindex{imagem}} de $f$.
O conjunto $B$ é chamado de \emph{\myindex{contradom\'inio}} de $f$.
\end{defn}

É poss\'ivel que a imagem $R(f)$ seja um subconjunto pr\'oprio do
contradom\'inio $B$, ao passo que o dom\'inio de $f$ é sempre igual a $A$.
Em geral, supomos que o dom\'inio de $f$ n\~ao seja vazio.

\begin{example}
Em Cálculo, você está mais familiarizado com funções de números reais em
números reais. No entanto, você também já viu outros tipos de funções.
A derivada é uma função que leva o conjunto das funções diferenciáveis ao
conjunto de todas as funções.
Outro exemplo é a transformada de Laplace, que também leva funções a funções.
Ainda outro exemplo é a função que recebe uma função contínua $g$ definida
no intervalo $[0,1]$ e retorna o número $\int_0^1 g(x) \,dx$.
\end{example}

\begin{defn}
Considere uma função $f \colon A \to B$. Definimos
a \emph{\myindex{imagem}} (ou \emph{\myindex{imagem direta}}) de um
subconjunto $C \subset A$ por
\glsadd{not:dirimage}
\begin{equation*}
f(C) \coloneqq \bigl\{ f(x) \in B : x \in C \bigr\} .
\end{equation*}
Definimos a \emph{\myindex{imagem inversa}} de um subconjunto
$D \subset B$ por
\glsadd{not:invimage}
\begin{equation*}
f^{-1}(D) \coloneqq \bigl\{ x \in A : f(x) \in D \bigr\} .
\end{equation*}
\end{defn}

Em particular, $R(f) = f(A)$: a imagem é a imagem direta do domínio $A$.

\begin{myfigureht}
\myincludepdft{funcimags_full}{%
Dois conjuntos de pontos estão identificados pelos números 1, 2, 3, 4 e
pelas letras a, b, c, d.
Há uma seta de 1 para b, de 2 para d, de 3 para c e de 4 para b. À direita,
uma lista apresenta exemplos de imagens:
A imagem, por f, do conjunto que contém 1, 2, 3 e 4 é o conjunto que contém
b, c e d.
A imagem, por f, do conjunto que contém 1, 2 e 4 é o conjunto que contém b
e d.
A imagem, por f, do conjunto que contém 1 é o conjunto que contém b.
A imagem inversa, por f, do conjunto que contém a, b e c é o conjunto que
contém 1, 3 e 4.
A imagem inversa, por f, do conjunto que contém a é o conjunto vazio.
A imagem inversa, por f, do conjunto que contém b é o conjunto que contém 1
e 4.}
\caption{Exemplo de imagens diretas e inversas para a fun\c{c}\~ao
$f \colon \{ 1,2,3,4 \} \to \{ a,b,c,d \}$ definida por
$f(1) \coloneqq b$,
$f(2) \coloneqq d$,
$f(3) \coloneqq c$,
$f(4) \coloneqq b$.\label{figfuncimags}}
\end{myfigureht}

\begin{example}
Defina a função $f \colon \R \to \R$ por
$f(x) \coloneqq \sin(\pi x)$. Então $f\bigl([0,\nicefrac{1}{2}]\bigr) = [0,1]$,
$f^{-1}\bigl(\{0\}\bigr) = \Z$, etc.
\end{example}

\begin{prop} \label{st:propinv}
Considere $f \colon A \to B$. Sejam $C, D$ subconjuntos de $B$. Então
\begin{align*}
& f^{-1}( C \cup D) = f^{-1} (C) \cup f^{-1} (D) , \\
& f^{-1}( C \cap D) = f^{-1} (C) \cap f^{-1} (D) , \\
& f^{-1}( C^c) = {\bigl( f^{-1} (C) \bigr)}^c .
\end{align*}
\end{prop}

Interprete a última linha da proposição como
$f^{-1}( B \setminus C) = A \setminus f^{-1} (C)$.

\begin{proof}
Começamos pela união. Se $x \in f^{-1}( C \cup D)$, então a função leva
$x$ a um elemento de $C$ ou de $D$, isto é, $f(x) \in C$ ou $f(x) \in D$.
Logo, $f^{-1}( C \cup D) \subset f^{-1} (C) \cup f^{-1} (D)$.
Reciprocamente, se $x \in f^{-1}(C)$, então $x \in f^{-1}(C \cup D)$.
O mesmo vale para $x \in f^{-1}(D)$. Portanto,
$f^{-1}( C \cup D) \supset f^{-1} (C) \cup f^{-1} (D)$, e obtemos a
igualdade.

O restante da demonstração fica como exercício.
\end{proof}

Para imagens diretas, o melhor que podemos afirmar é o resultado mais fraco
a seguir.

\begin{prop} \label{st:propfor}
Considere $f \colon A \to B$. Sejam $C, D$ subconjuntos de $A$. Então
\begin{align*}
& f( C \cup D) = f (C) \cup f (D) , \\
& f( C \cap D) \subset f (C) \cap f (D) .
\end{align*}
\end{prop}

Deixamos a demonstração como exercício.

\begin{defn}
Seja $f \colon A \to B$ uma função.
Dizemos que $f$ é
\emph{\myindex{injetiva}} ou
\emph{\myindex{um a um}} se $f(x_1) = f(x_2)$ implica $x_1 = x_2$. Em
outras palavras, $f$ é injetiva se, para todo $y \in B$, o conjunto
$f^{-1}(\{y\})$ é vazio ou contém um único elemento.
Chamamos tal $f$ de uma \emph{\myindex{inje\c{c}\~ao}}.

Se $f(A) = B$, dizemos que $f$ é
\emph{\myindex{sobrejetiva}} ou
\emph{\myindex{sobrejetora}}.
Em outras palavras, $f$ é sobrejetiva se a imagem e o contradomínio de $f$
são iguais.
Chamamos tal $f$ de uma \emph{\myindex{sobreje\c{c}\~ao}}.

Se $f$ é ao mesmo tempo sobrejetiva e injetiva, dizemos que $f$ é
\emph{\myindex{bijetiva}} ou que $f$ é uma
\emph{\myindex{bije\c{c}\~ao}}.
\end{defn}

Quando $f \colon A \to B$ é uma bijeção, a imagem inversa de um elemento,
$f^{-1}(\{y\})$, é sempre um único elemento de $A$. Passamos então a
considerar $f^{-1}$ como uma função
\glsadd{not:invfunction}
$f^{-1} \colon B \to A$ e escrevemos simplesmente $f^{-1}(y)$.
Nesse caso, chamamos $f^{-1}$ de \emph{\myindex{fun\c{c}\~ao inversa}} de $f$.
Por exemplo, para a bijeção $f \colon \R \to \R$ definida por
$f(x) \coloneqq x^3$, temos
$f^{-1}(x) = \sqrt[3]{x}$.

%A final piece of notation for functions that
%we need is the \emph{\myindex{composition of functions}}.

\begin{defn}
Considere $f \colon A \to B$ e $g \colon B \to C$. A
\emph{composi\c{c}\~ao}\index{composi\c{c}\~ao de fun\c{c}\~oes}
das funções $f$ e $g$ é a função
$g \circ f \colon A \to C$ definida por
\glsadd{not:composition}
\begin{equation*}
(g \circ f)(x) \coloneqq g\bigl(f(x)\bigr) .
\end{equation*}
\end{defn}

Por exemplo, se $f \colon \R \to \R$ é dada por $f(x)\coloneqq x^3$ e
$g \colon \R \to \R$ é dada por $g(y) = \sin(y)$, então
$(g \circ f)(x) = \sin(x^3)$.
Deixamos como exercício simples para o leitor mostrar que a composição de
aplicações injetivas é injetiva e a composição de aplicações sobrejetivas é
sobrejetiva. Portanto, a composição de bijeções é uma bijeção.

\subsection{Rela\c{c}\~oes e classes de equival\^encia}

Muitas vezes, comparamos dois objetos de alguma maneira. Por exemplo,
dizemos $1 < 2$ para números naturais, $\nicefrac{1}{2} = \nicefrac{2}{4}$
para números racionais ou $\{ a,c \} \subset \{ a,b,c \}$ para conjuntos.
Os símbolos `$<$', `$=$' e `$\subset$' são exemplos de relações.

\begin{defn}
Seja $A$ um conjunto. Uma \emph{\myindex{rela\c{c}\~ao bin\'aria}}\index{rela\c{c}\~ao}
em $A$ é um subconjunto $\sR \subset A \times A$,
formado pelos pares para os quais dizemos que a relação vale.
Em vez de escrever $(a,b) \in \sR$, escrevemos
$a \, \sR \, b$.
\end{defn}

\begin{example}
Considere $A \coloneqq \{ 1,2,3 \}$.

Considere a relação `$<$'. O conjunto correspondente de pares é
$\bigl\{ (1,2), (1,3), (2,3) \bigr\}$. Assim, $1 < 2$ vale, pois $(1,2)$
pertence ao conjunto de pares correspondente, mas $3 < 1$ não vale, pois
$(3,1)$ não pertence ao conjunto.

De modo semelhante, a relação `$=$' é definida pelo conjunto de pares
$\bigl\{ (1,1), (2,2), (3,3) \big\}$.

Qualquer subconjunto de $A \times A$ é uma relação. Se definirmos a relação
$\dagger$ pelo conjunto $\bigl\{ (1,2), (2,1), (2,3), (3,1) \bigr\}$, então
As afirmações $1 \dagger 2$ e $3 \dagger 1$ são verdadeiras, mas
$1 \dagger 3$ não é.
\end{example}

\begin{defn}
%FIXMEevillayouthack
\pagebreak[2]
A relação
$\sR$ em um conjunto $A$ é dita
\begin{enumerate}[(i)]
\item
\emph{Reflexiva}\index{rela\c{c}\~ao reflexiva} se $a\, \sR \, a$ para
todo $a \in A$.
\item
\emph{Simétrica}\index{rela\c{c}\~ao sim\'etrica} se $a\, \sR \, b$ implica
$b \, \sR \, a$.
\item
\emph{Transitiva}\index{rela\c{c}\~ao transitiva} se $a\, \sR \, b$ e
$b \, \sR \, c$ implicam $a\, \sR \, c$.
\end{enumerate}
Se $\sR$ é reflexiva, simétrica e transitiva, dizemos que é uma
\emph{\myindex{rela\c{c}\~ao de equival\^encia}}.
\end{defn}

\begin{example}
Seja $A \coloneqq \{ 1,2,3 \}$.
A relação `$<$' é transitiva, mas não é reflexiva nem simétrica. A
relação `$\leq$', definida por
$\bigl\{ (1,1), (1,2), (1,3), (2,2), (2,3), (3,3) \bigr\}$
é reflexiva e transitiva, mas não é simétrica.
Por fim, a relação `$\star$', definida por
$\bigl\{ (1,1), (1,2), \allowbreak (2,1), \allowbreak (2,2), \allowbreak (3,3) \bigr\}$
é uma relação de equivalência.
\end{example}

As relações de equivalência são úteis porque dividem um conjunto em
subconjuntos de elementos \myquote{equivalentes}.

\begin{defn}
Seja $A$ um conjunto e $\sR$ uma relação de equivalência.
Uma \emph{\myindex{classe de equival\^encia}} de $a \in A$, geralmente
denotada por $[a]$,\glsadd{not:equivclass} é o conjunto
$\{ x \in A : a\, \sR \, x \}$.
\end{defn}

Por exemplo, para a relação `$\star$' acima, há duas classes de equivalência:
$[1] = [2] = \{ 1,2 \}$ e $[3] = \{ 3 \}$.

A reflexividade garante que $a \in [a]$. A simetria garante que, se
$b \in [a]$, então $a \in [b]$. Por fim, a transitividade garante que,
se $b \in [a]$ e $c \in [b]$, então $c \in [a]$.
Em particular, temos a seguinte proposição, cuja demonstração fica como
exercício.

\begin{prop} \label{prop:equivclasses}
Se $\sR$ é uma relação de equivalência em um conjunto $A$,
então cada $a \in A$ pertence a exatamente uma classe de equivalência.
Além disso, $a\, \sR \, b$ se, e somente se, $[a] = [b]$.
\end{prop}

\begin{example} \label{example:ratnums}
O conjunto dos números racionais pode ser definido como o conjunto das
classes de equivalência de pares ordenados formados por um número inteiro e
um número natural, isto é, de elementos de $\Z \times \N$. A relação é definida por
$(a,b) \sim (c,d)$ sempre que $ad = bc$.
Deixamos como exercício provar que `$\sim$' é uma relação de equivalência.
Em geral, a classe de equivalência $\bigl[(a,b)\bigr]$ é denotada por
$\nicefrac{a}{b}$.
\end{example}

\subsection{Cardinalidade}

Uma questão sutil, mas fundamental, na teoria dos conjuntos, e que causa
muita confusão entre estudantes iniciantes, é a cardinalidade, ou o
\myquote{tamanho} dos conjuntos.
%The concept of cardinality is fundamental in modern mathematics in general and
%in analysis in particular.
De fato, nesta seção, veremos o primeiro teorema realmente inesperado.

\begin{defn}
Sejam $A$ e $B$ conjuntos. Dizemos que $A$ e $B$ têm a mesma
\emph{\myindex{cardinalidade}}
quando existe uma bijeção $f \colon A \to B$. Denotamos
por $\sabs{A}$ \glsadd{not:cardinality} a classe de equivalência de todos os
conjuntos com a mesma cardinalidade que $A$ e chamamos $\sabs{A}$ simplesmente
de cardinalidade de $A$.
\end{defn}

Por exemplo, $\{ 1,2,3 \}$ tem a mesma cardinalidade que $\{ a,b,c \}$ se
definirmos a bijeção $f(1) \coloneqq a$, $f(2) \coloneqq b$,
$f(3) \coloneqq c$.
Evidentemente, essa bijeção não é única.

A relação \myquote{ter a mesma cardinalidade} é, de fato, uma relação de
equivalência.
A função identidade, $f(x) \coloneqq x$, é uma bijeção que mostra a
reflexividade.
Se $f$ é uma bijeção, então $f^{-1}$ também é, o que mostra a simetria.
Se $f \colon A \to B$ e $g \colon B \to C$ são bijeções, então
$g \circ f$ é uma bijeção entre $A$ e $C$, o que mostra a transitividade.
Um conjunto $A$ tem a mesma cardinalidade que o conjunto vazio se, e somente
se $A$ é o próprio conjunto vazio: se $B$ não é vazio, não existe função
$f \colon B \to \emptyset$.
Em particular, não existe bijeção entre $B$ e $\emptyset$.

\begin{defn}
Se $A$ tem a mesma cardinalidade que $\{ 1,2,3,\ldots,n \}$
para algum $n \in \N$,
escrevemos $\sabs{A} \coloneqq n$. Se $A$ é vazio, escrevemos
$\sabs{A} \coloneqq 0$.
Em ambos os casos, dizemos que $A$ é
\emph{\myindex{finito}}.
Dizemos que $A$ é \emph{\myindex{infinito}} ou tem
\myquote{cardinalidade infinita} se $A$ não for finito.
\end{defn}

Deixamos como exercício verificar que a notação $\sabs{A} = n$ é justificada.
Isto é, para cada conjunto finito não vazio $A$, existe um único número
natural $n$ tal que há uma bijeção de $A$ em $\{ 1,2,3,\ldots,n \}$.

Podemos ordenar conjuntos pelo tamanho (cardinalidade).

\begin{defn} \label{def:comparecards}
Escrevemos
\begin{equation*}
\sabs{A} \leq \sabs{B}
\end{equation*}
se existe uma injeção de $A$ em $B$. Escrevemos $\sabs{A} = \sabs{B}$
se $A$ e $B$ têm a mesma cardinalidade. Escrevemos $\sabs{A} < \sabs{B}$
se $\sabs{A} \leq \sabs{B}$, mas $A$ e $B$ não têm a mesma cardinalidade.
\end{defn}

Enunciamos, sem demonstração, que
$A$ e $B$ têm a mesma cardinalidade se, e somente se,
$\sabs{A} \leq \sabs{B}$ e
$\sabs{B} \leq \sabs{A}$. Esse resultado é o teorema de
\myindex{Cantor--Bernstein--Schr\"oder}.
Além disso, para quaisquer dois conjuntos $A$ e $B$,
sempre podemos escrever $\sabs{A} \leq \sabs{B}$ ou
$\sabs{B} \leq \sabs{A}$. As questões relacionadas a essa
última afirmação são muito sutis. Como não precisamos de nenhum
desses dois resultados, omitimos as demonstrações.

Os casos realmente interessantes de cardinalidade são os de conjuntos
infinitos.
Vamos distinguir dois tipos de cardinalidade infinita.

\begin{defn}
Se $\sabs{A} = \sabs{\N}$, dizemos que $A$ é
\emph{\myindex{infinito enumer\'avel}}.
Se $A$ é finito ou infinito enumerável, dizemos que $A$
é \emph{\myindex{enumer\'avel}}.
Se $A$ não é enumerável,
dizemos que $A$ é \emph{\myindex{n\~ao enumer\'avel}}.
\end{defn}

A cardinalidade de $\N$ costuma ser denotada por
$\aleph_0$ (lê-se alef zero)\footnote{Para os fãs do programa de TV
\emph{Futurama}, há uma sala de cinema em um episódio
chamada de $\aleph_0$-plex.}.

\begin{example}
O conjunto dos números naturais pares tem a mesma cardinalidade que $\N$.
Demonstração: Seja $E \subset \N$ o conjunto dos números naturais pares.
Dado $k \in E$, escrevemos $k=2n$ para algum $n \in \N$.
Então $f(n) \coloneqq 2n$ define uma bijeção $f \colon \N \to E$.
\end{example}

Mencionamos, sem demonstração, a seguinte caracterização dos conjuntos
infinitos: \emph{Um conjunto é infinito se, e somente se, está em
correspondência biunívoca com algum subconjunto próprio de si mesmo}.

\begin{example}
$\N \times \N$ é um conjunto infinito enumerável. Demonstração: Organizemos
os elementos de $\N \times \N$ da seguinte maneira:
$(1,1)$, $(1,2)$, $(2,1)$, $(1,3)$, $(2,2)$, $(3,1)$, \ldots. Isto é,
primeiro listamos todos os elementos cujas duas entradas somam $k$,
depois listamos todos aqueles cujas entradas somam $k+1$, e assim por diante.
Definimos uma bijeção com $\N$ fazendo 1 corresponder a $(1,1)$,
2 corresponder a $(1,2)$, e assim por diante. Veja \figureref{fig:NcrossNcard}.
\begin{myfigureht}
\myincludepdft{NcrossNcard}{%
Uma grade de pares (linha, coluna) é percorrida por uma seta que vai
de (1,1) a (1,2), depois a
(2,1), diagonalmente para baixo e para a esquerda, e então a (1,3); em seguida,
diagonalmente para baixo e para a esquerda até (2,2) e depois até (3,1). Uma
seta segue até (1,4) e continua diagonalmente para baixo e para a esquerda por
(2,3), (3,2) e (4,1). Reticências indicam que o diagrama é infinito.}
\caption{Mostrando que $\N \times \N$ é enumerável.\label{fig:NcrossNcard}}
\end{myfigureht}
\end{example}

\begin{example}
O conjunto dos números racionais é enumerável. Demonstração (informal):
Para os racionais positivos, seguimos o mesmo procedimento do exemplo
anterior, escrevendo
$\nicefrac{1}{1}$, $\nicefrac{1}{2}$, $\nicefrac{2}{1}$ etc. No entanto,
omitimos as frações (como $\nicefrac{2}{2}$)
que já apareceram. A lista continua assim:
$\nicefrac{1}{3}$, $\nicefrac{3}{1}$, $\nicefrac{1}{4}$,
$\nicefrac{2}{3}$ etc.
Para listar todos os racionais, incluímos o $0$ e os números negativos:
$0$,
$\nicefrac{1}{1}$,
$\nicefrac{-1}{1}$,
$\nicefrac{1}{2}$,
$\nicefrac{-1}{2}$ etc.
\end{example}

Para completar, mencionamos as seguintes afirmações que aparecem nos
exercícios.
\emph{Se $A \subset
B$ e $B$ é enumerável, então $A$ é enumerável.} A contrapositiva dessa
afirmação diz que, se $A$ não é enumerável, então $B$ também não é
enumerável.
Como consequência, se $\sabs{A} < \sabs{\N}$, então $A$ é
finito.
De modo semelhante, se $B$ é finito e $A \subset B$, então $A$ é finito.

%FIXMEevillayouthack
\pagebreak[2]
Apresentamos agora o primeiro resultado realmente surpreendente sobre
cardinalidade.
Para isso, precisamos de uma notação para
o conjunto de todos os subconjuntos de um conjunto.

\begin{defn}
O \emph{\myindex{conjunto das partes}} de um conjunto $A$, denotado por
$\sP(A)$,\glsadd{not:powerset} é o conjunto de todos os subconjuntos de $A$.
\end{defn}

Por exemplo, se $A \coloneqq \{ 1,2\}$,
então $\sP(A) = \bigl\{ \emptyset, \{ 1 \}, \{ 2 \}, \{ 1, 2 \} \bigr\}$.
Em particular, $\sabs{A} = 2$ e $\babs{\sP(A)} = 4 = 2^2$.
Em geral,
para um conjunto finito $A$ de cardinalidade $n$,
a cardinalidade de $\sP(A)$ é $2^n$.
Deixamos esse fato como exercício.
Portanto, para um conjunto finito $A$,
a cardinalidade de $\sP(A)$ é estritamente
maior que a
cardinalidade de $A$. O fato inesperado e
surpreendente é que essa afirmação também vale para conjuntos infinitos.

\begin{thm}[Cantor%
\footnote{Em homenagem ao matemático alemão
\href{https://en.wikipedia.org/wiki/Georg_Cantor}{Georg Ferdinand Ludwig
Philipp Cantor} (1845--1918).}]
\index{teorema de Cantor}
\label{cantorspowersetthm}
Seja $A$ um conjunto. Então
$\sabs{A} < \babs{\sP(A)}$.
Em particular, não existe função sobrejetora de
$A$ em $\sP(A)$.
\end{thm}

\begin{proof}
Existe uma injeção $f \colon A \to \sP(A)$:
para $x \in A$, seja $f(x) \coloneqq \{ x \}$. Assim,
$\sabs{A} \leq \babs{\sP(A)}$.

Para concluir a demonstração, devemos mostrar que
nenhuma função $g \colon A \to \sP(A)$ é sobrejetora.
Suponha que
$g \colon A \to \sP(A)$ seja uma função. Então, para $x \in A$,
$g(x)$ é um subconjunto de $A$. Definimos o conjunto
\begin{equation*}
B \coloneqq \bigl\{ x \in A : x \notin g(x) \bigr\} .
\end{equation*}
Afirmamos que $B$ não pertence à imagem de $g$ e, portanto, $g$ não é
sobrejetora. Suponha, por contradição, que exista um $x_0$ tal que
$g(x_0) = B$.
Ou $x_0 \in B$ ou $x_0 \notin B$. Se $x_0 \in B$, então $x_0 \notin
g(x_0) = B$, o que é uma contradição. Se $x_0 \notin B$, então $x_0 \in
g(x_0) = B$, o que também é uma contradição. Logo, tal $x_0$ não
existe. Portanto, $B$ não pertence à imagem de $g$, e $g$ não é
sobrejetora. Como $g$ era uma função arbitrária, não existe nenhuma função
sobrejetora.
\end{proof}

Uma consequência particular deste teorema é a existência de conjuntos não
enumeráveis, pois $\sP(\N)$ necessariamente não é enumerável.
Um fato relacionado é que
o conjunto dos números reais (que estudaremos no próximo capítulo) não é
enumerável.
A existência de conjuntos não enumeráveis pode parecer contraintuitiva, e o
teorema causou bastante controvérsia na época
em que foi anunciado. O teorema não afirma apenas que existem conjuntos não
enumeráveis: afirma também que, de fato, existem conjuntos infinitos cada vez
maiores,
$\N$, $\sP(\N)$,
$\sP(\sP(\N))$, $\sP(\sP(\sP(\N)))$ etc.


\subsection{Exercícios}

\begin{exercise}
Mostre
$A \setminus (B \cap C) = (A \setminus B) \cup (A \setminus C)$.
\end{exercise}

\begin{exercise}
Demonstre que o princípio da indução forte é equivalente ao princípio da
indução usual.
\end{exercise}

\begin{exercise}
Complete a demonstração de \propref{st:propinv}.
\end{exercise}

\begin{exercise}
\leavevmode
\begin{enumerate}[a)]
\item
Demonstre \propref{st:propfor}.
\item
Encontre um exemplo em que falha a igualdade de conjuntos
em
$f( C \cap D) \subset f (C) \cap f (D)$.
Isto é, encontre $f$, $A$, $B$, $C$ e $D$ tais que
$f( C \cap D)$ seja um subconjunto próprio de $f(C) \cap f(D)$.
\end{enumerate}
\end{exercise}

\begin{exercise}[Difícil]
Demonstre que, se $A$ é finito e não vazio, então existe um único
$n \in \N$ tal que há uma bijeção entre $A$ e $\{ 1, 2, 3, \ldots, n \}$.
Em outras palavras, a notação $\sabs{A} \coloneqq n$ é justificada.
Dica: mostre que, se $n > m$, então não existe injeção de
$\{ 1, 2, 3, \ldots, n \}$ em
$\{ 1, 2, 3, \ldots, m \}$.
\end{exercise}


\begin{exercise}
Demonstre:
\begin{enumerate}[a)]
\item $A \cap (B \cup C) = (A \cap B) \cup (A \cap C)$.
\item $A \cup (B \cap C) = (A \cup B) \cap (A \cup C)$.
\end{enumerate}
\end{exercise}

\begin{samepage}
\begin{exercise}
Denotemos por $A \Delta B$ a
\emph{\myindex{diferen\c{c}a sim\'etrica}}, isto é, o conjunto de todos os elementos que
pertencem a $A$ ou a $B$, mas não a ambos os conjuntos, $A$ e $B$.
\begin{enumerate}[a)]
\item
Desenhe um diagrama de Venn para
$A \Delta B$.
\item
Mostre que $A \Delta B = (A \setminus B) \cup (B \setminus A)$.
\item
Mostre que $A \Delta B = (A \cup B) \setminus ( A \cap B)$.
\end{enumerate}
\end{exercise}
\end{samepage}

\begin{exercise}
Para cada $n \in \N$, seja $A_n \coloneqq \{ (n+1)k : k \in \N \}$.
\begin{enumerate}[a)]
\item Determine $A_1 \cap A_2$.
\item Determine $\bigcup_{n=1}^\infty A_n$.
\item Determine $\bigcap_{n=1}^\infty A_n$.
\end{enumerate}
\end{exercise}

\begin{samepage}
\begin{exercise}
Determine $\sP(S)$ (o conjunto das partes) para cada um dos conjuntos a seguir:
\begin{enumerate}[a)]
\item $S = \emptyset$,
\item $S = \{1\}$,
\item $S = \{1,2\}$,
\item $S = \{1,2,3,4\}$.
\end{enumerate}
\end{exercise}
\end{samepage}


\begin{exercise}
Sejam $f \colon A \to B$ e $g \colon B \to C$ funções.
\begin{enumerate}[a)]
\item
Demonstre que, se $g \circ f$ é injetiva, então $f$ é injetiva.
\item
Demonstre que, se $g \circ f$ é sobrejetiva, então $g$ é sobrejetiva.
\item
Encontre um exemplo explícito em que $g \circ f$ é bijetiva, mas nem $f$
nem $g$ é bijetiva.
\end{enumerate}
\end{exercise}

\begin{exercise}
\pagebreak[2]
Demonstre por indução que $n < 2^n$ para todo $n \in \N$.
\end{exercise}

\begin{exercise}
Mostre que, para um conjunto finito $A$ de cardinalidade $n$, a cardinalidade
de $\sP(A)$ é $2^n$.
\end{exercise}

\begin{exercise}
Demonstre que $\frac{1}{1\cdot 2} +
\frac{1}{2\cdot 3} + \cdots + \frac{1}{n(n+1)} = \frac{n}{n+1}$
para todo $n \in \N$.
\end{exercise}

\begin{exercise}
Demonstre que $1^3 + 2^3 + \cdots + n^3 = {\left( \frac{n(n+1)}{2} \right)}^2$
para todo $n \in \N$.
\end{exercise}

\begin{exercise}
Demonstre que $n^3 + 5n$ é divisível por $6$ para todo $n \in \N$.
\end{exercise}

\begin{exercise}
Encontre o menor $n \in \N$ tal que $2{(n+5)}^2 < n^3$ e chame-o de $n_0$.
Mostre que $2{(n+5)}^2 < n^3$ para todo $n \geq n_0$.
\end{exercise}

\begin{exercise}
Encontre todos os $n \in \N$ tais que $n^2 < 2^n$.
\end{exercise}

\begin{exercise}
Demonstre a
\hyperlink{wop:link}{propriedade da boa ordenação} de $\N$ usando o
\hyperref[induction:thm]{princípio da indução}.
\end{exercise}


\begin{exercise}
Dê um exemplo de uma coleção infinita enumerável de conjuntos finitos
$A_1, A_2, \ldots$ cuja união não seja um conjunto finito.
\end{exercise}

\begin{exercise}
Dê um exemplo de uma coleção infinita enumerável de conjuntos infinitos
$A_1, A_2, \ldots$ em que $A_j \cap A_k$ seja infinito para quaisquer $j$ e $k$,
mas
$\bigcap_{j=1}^\infty A_j$
seja não vazio e finito.
\end{exercise}

\begin{exercise}
Suponha que $A \subset B$ e que $B$ seja finito. Demonstre que $A$ é finito.
Isto é, se $A$ não é vazio, construa uma bijeção de $A$ em $\{ 1,2,\ldots,n \}$.
\end{exercise}

\begin{exercise}
Demonstre \propref{prop:equivclasses}. Isto é, demonstre que, se $\sR$ é uma
relação de equivalência em um conjunto $A$, então cada $a \in A$ pertence a
exatamente uma classe de equivalência. Em seguida, demonstre que
$a \, \sR \, b$ se, e somente se, $[a] =
[b]$.
\end{exercise}

\begin{exercise}
Demonstre que a relação $\sim$ em \exampleref{example:ratnums} é uma relação
de equivalência.
\end{exercise}

\begin{exercise}
\leavevmode
\begin{enumerate}[a)]
\item
Suponha que $A \subset B$ e que $B$ seja infinito enumerável. Construindo
uma bijeção, mostre que $A$ é
enumerável (isto é, $A$ é vazio, finito ou infinito enumerável).
\item
Use o item (a) para mostrar que, se $\sabs{A} < \sabs{\N}$, então $A$ é finito.
\end{enumerate}
\end{exercise}

\begin{exercise}[Desafiador] \label{exercise:countsubsetbij}
Suponha que $\sabs{\N} \leq \sabs{S}$, ou, em outras palavras, que $S$
contenha um subconjunto infinito enumerável.
Mostre que existe um subconjunto infinito enumerável $A \subset S$
tal que há uma bijeção entre $S \setminus A$ e $S$.
\end{exercise}

\begin{exercise}
Demonstre as versões infinitas das leis de De Morgan.
Suponha que $A$ seja um conjunto e que $B_\lambda$ seja uma coleção de
conjuntos para $\lambda \in I$. Demonstre que
\begin{equation*}
A \setminus \biggl(
\bigcup_{\lambda \in I} B_\lambda
\biggr)
=
\bigcap_{\lambda \in I}
(
A \setminus
B_\lambda
) ,
\qquad
A \setminus \biggl(
\bigcap_{\lambda \in I} B_\lambda
\biggr)
=
\bigcup_{\lambda \in I}
(
A \setminus
B_\lambda
) .
\end{equation*}
\end{exercise}

\begin{exercise}
Suponha que $f \colon A \to B$ seja uma função e que, para $\lambda \in I$,
tenhamos uma coleção de subconjuntos $C_\lambda \subset A$ e
$D_\lambda \subset B$. Demonstre que
\begin{equation*}
f^{-1} \biggl(
\bigcup_{\lambda \in I} D_\lambda
\biggr)
=
\bigcup_{\lambda \in I}
f^{-1}(D_\lambda) ,
\quad
f^{-1} \biggl(
\bigcap_{\lambda \in I} D_\lambda
\biggr)
=
\bigcap_{\lambda \in I}
f^{-1}(D_\lambda) ,
\end{equation*}
e
\begin{equation*}
f \biggl(
\bigcup_{\lambda \in I} C_\lambda
\biggr)
=
\bigcup_{\lambda \in I}
f(C_\lambda) ,
\quad
f \biggl(
\bigcap_{\lambda \in I} C_\lambda
\biggr)
\subset
\bigcap_{\lambda \in I}
f(C_\lambda) .
\end{equation*}
\end{exercise}

